%% Facilities, Equipment, and Other Resources
%% NSF 26-506 (PESOSE), Track 3 -- RedTwin
%%
%% Compile:  pdflatex facilities.tex   (or `make facilities` from the repo root)
%%
%% COMPLIANCE NOTES (PAPPG 24-1):
%%   - II.D.2.g: narrative in nature; describe ONLY directly applicable resources;
%%     cover the resources that the organization, its COLLABORATORS, and its
%%     SUBRECIPIENTS will provide; and include NO quantifiable financial
%%     information.  (The previous draft's "$800 million in annual research
%%     expenditures" violated the last of these.)
%%   - II.D.2.d(iv): unfunded collaborations are documented by a letter AND
%%     described here.  Section 3 below is that description.
%%   - II.C.2: the font, spacing, and margin rules apply to supplementary
%%     documents too.  Times New Roman 11pt, 1in margins, no page numbers
%%     (Research.gov paginates).
%%   - URLs ARE permitted here; the no-URL rule is scoped to the Project
%%     Description only.
%%   - NSF expects resources named in this section to be provided if the award
%%     is made, so collaborator commitments are stated conditionally, in the
%%     terms their letters actually use.

\documentclass[11pt,letterpaper]{article}
\usepackage[letterpaper,margin=1in,heightrounded,
            headheight=0pt,headsep=0pt,footskip=0pt]{geometry}
\usepackage[T1]{fontenc}
\usepackage[utf8]{inputenc}
\usepackage{times}
\usepackage{microtype}
\usepackage{titlesec}
\usepackage{enumitem}
\usepackage{xcolor}
\usepackage{xspace}

\newcommand{\sysname}{RedTwin\xspace}

%% Conspicuous placeholder for facts only the PIs can confirm.
\newcommand{\todoverify}[1]{%
  \textbf{\textcolor{red}{[PI TO SUPPLY: #1]}}}

\setlength{\parindent}{0pt}
\setlength{\parskip}{0.55em}
\titleformat{\section}{\normalfont\large\bfseries}{\thesection.}{0.5em}{}
\titlespacing*{\section}{0pt}{1.3ex plus .2ex}{0.5ex}
\titleformat{\subsection}[runin]{\normalfont\bfseries\itshape}{}{0em}{}[.]
\titlespacing*{\subsection}{0pt}{0.6ex plus .2ex}{0.5em}
\pagestyle{empty}

\begin{document}

\begin{center}
  {\large\bfseries Facilities, Equipment, and Other Resources}
\end{center}

\vspace{0.3em}

This section describes the resources that the University of California, Santa
Barbara (the lead organization), Arizona State University (the subrecipient),
and the project's collaborating organizations will make available to
\sysname. Only resources that are directly applicable to the proposed work are
described. The organizations named in Section~3 receive no funds under this
award.

\section{University of California, Santa Barbara (lead organization)}

UC Santa Barbara is a member of the Association of American Universities and a
Carnegie R1 institution. Its College of Engineering and Department of Computer
Science provide office and laboratory space, computing, and administrative
support to the PIs and their graduate researchers.

\subsection{Laboratory space}
PIs Vigna and Kruegel direct the Security Laboratory (SecLab), located on the
second floor of Harold Frank Hall. SecLab provides desk and meeting space for
up to twenty researchers, and it includes a machine room with the rack space,
power, and cooling required by the server clusters this project will use.

\subsection{Computing for twin construction and analysis}
The proposed work instantiates many virtualized hosts at once, since each
digital twin reproduces a deployment of roughly twenty to one hundred
components, and it sustains fuzzing and large-language-model inference against
those twins on a nightly and weekly cadence. Three tiers of resource cover
this. SecLab operates its own Linux servers and workstations for day-to-day
development and for the smaller reference networks. The Department of Computer
Science maintains a high-performance computing server of four nodes carrying
eight NVIDIA A100 GPUs each, which supports model inference and the
component-level analyses of Thrust~3. Beyond the department, the Center for
%% Specifications below follow the CSC "Get Started" page,
%% https://csc.cnsi.ucsb.edu/get-started (checked August 2026). That page lists
%% the V100 nodes as V100/40GB; the memory figure is omitted here because the
%% V100 normally ships in 16GB and 32GB variants. Confirm before submission if
%% you want the figure included. Condo clusters (Braid, Braid2) are excluded
%% because access depends on node ownership. Per-node condo pricing is on that
%% page and must NOT be quoted here: PAPPG II.D.2.g forbids quantifiable
%% financial information in this section.
Scientific Computing (CSC) at the California NanoSystems Institute operates
clusters that are open to any researcher on campus. The Pod cluster provides
40-core CPU nodes for general-purpose computing together with fifteen
quad-GPU NVIDIA V100 nodes connected by SXM and NVLINK, and the Knot7 cluster
serves large parameter sweeps and jobs that exceed a local workstation. Both
are available at no charge to campus researchers, subject to acknowledgment in
resulting publications, and CSC staff hold backgrounds in science or
mathematics and assist users with porting, developing, and optimizing code.
The scale-out runs of Section~7 will use Pod to instantiate the larger twins
and to execute the planner against them, while its GPU nodes carry the
inference load of the agent teams alongside the departmental A100 server. The
UCSB campus cloud, part of the NSF Aristotle project, provides additional
elastic capacity.

\subsection{Isolated environment for digital twins}
Section~6 of the Project Description commits to holding every twin on an
isolated network, restricting access to the project team and the owning
organization, and destroying twins at the end of an engagement. SecLab
provides the segregated environment in which that commitment is met.
The SecLab manages an isolated network that has a reserved IP space separate from the UCSB network. 
The SecLab administrators have a long experience in running experiments that need to be contained in a contained environment (from malware analysis to the reproduction of vulnerabilities).

\subsection{Reference infrastructure}
The Global Accessible Test Infrastructure (GATE), developed and maintained
under the ACTION AI Institute, provides emulated topologies representing smart
cities, chemical plants, water systems, and manufacturing plants, and is
available to the research community. GATE topologies will serve as input
reference infrastructures for the evaluation of Section~7. No funds are
requested for them under this award.

\subsection{Personnel provided through the Early Research Scholars Program}
The Early Research Scholars Program (ERSP) is a year-long, dual-mentored
research apprenticeship in the Department of Computer Science that pairs small
undergraduate teams with a faculty PI and a graduate mentor. Its director,
Ziad Matni, has committed to placing a dedicated team of three to five
second-year undergraduate researchers with PIs Vigna and Kruegel, co-advised
by an assigned SecLab graduate mentor. These students will contribute to
curating candidate open-source infrastructure configurations, to designing test
cases that exercise the logic-based invariant checks and the attack-graph
preservation analysis of Thrust~2, to the empirical evaluation of the
vulnerability-chaining agents, and to the design of the twinCTF challenges
described in Section~11. This effort is provided at no cost to the award.

\subsection{Campus systems for the operational evaluation}
Sections~6 and~7 describe the evaluation of \sysname on internal research
networks at UCSB, performed only under written authorization that specifies
the scope, the permitted observation techniques, and the rules of engagement.
This evaluation will be performed in collaboration with the UCSB Information Technology group, and specifically with the Campus Cybersecurity Services group.

\section{Arizona State University (subrecipient)}

Arizona State University is a member of the Association of American
Universities and a Carnegie R1 institution with very high research activity.

\subsection{Center and laboratory}
The Center for Cybersecurity and Trusted Foundations (CTF), part of ASU's
Advanced Capabilities for National Security Institute, is a multidisciplinary
center whose research spans binary analysis, vulnerability analysis, malware
attribution, automated binary analysis systems, and the collaboration of human
analysts with automated cyber reasoning systems. Its projects are driven by
real-world challenges and result in open-source prototypes. PI Doup\'e directs
the Center and co-directs, with Dr.~Shoshitaishvili, the Laboratory of
Security Engineering for Future Computing (SEFCOM), which occupies a
designated experimental research area within it. Beyond the hardware below,
CTF supplies the interdisciplinary connections through which this project can
reach collaborators in adjacent areas.

\subsection{Equipment applicable to this project}
SEFCOM operates approximately forty high-end servers and thirty PC systems
running Windows and UNIX variants. These provide the ASU-side capacity for the
component-level analysis of Thrust~3 and for the mock synthesis of Thrust~2.
SEFCOM also holds an ARM DSTREAM high-performance debug and trace unit, which
supports the firmware re-hosting that ASU contributes to Thrust~2 for
components that cannot be replicated directly.

\subsection{Educational infrastructure}
PI Doup\'e is a lead developer of pwn.college, and PIs Doup\'e and Bao are
Associate Directors of ASU's American Cybersecurity Education Institute, which
operates the publicly available pwn.college instance and facilitates classes
for colleges across the United States. This platform will host the open-source
infrastructure-security modules described in Section~11, giving them the same
distribution as the platform itself.

\subsection{Campus systems for the operational evaluation}
ASU's CTF has excellent relationships with ASU's University Technology Office (UTO).
These relationships will enable us to create digital twins of production, enterprise-level applications and deployments.
In addition, CTF will work closely with ASU's Chief Information Security Officer so that all discovered vulnerabilities are addressed in the production systems. 

\section{Resources provided by collaborating organizations}

The organizations in this section are not included in the budget and will
receive no funds under this award. Each has provided a letter of
collaboration, and the resources described below are those the letters
commit to, on the conditions those letters state.

\subsection{Sage Bionetworks}
Sage Bionetworks develops and maintains Synapse, an open-source repository and
collaboration environment for biomedical research data released under the
Apache License~2.0. Sage holds data resources of its own, acts as a steward
and broker for resources contributed by other institutions and consortia, and
maintains the governance machinery, including data-use agreements, consent
frameworks, access review, and de-identification practice, that determines
what may be shared and under what conditions. For the synthetic-data work of
Thrust~2, Sage has indicated that where a specific use case is well defined
and has been validated against the applicable consent terms, data-use
agreements, and access-review process, it expects to work with the team on
access to material suitable for deriving realistic, non-identifying
populations for hospital digital twins, meaning structure, distributions, and
relationships rather than records. Where direct access is not appropriate,
Sage expects to help specify what a defensible synthetic substitute would have
to preserve. Any such arrangement is subject to Sage's governance review and
to the conditions set by the original data contributors.

\subsection{The ZMap project}
Zakir Durumeri\'c, an original author and continuing maintainer of ZMap, has
indicated that the extensions developed under this project, including probe
modules, structured output formats, incremental-scan support, and constraint
handling, would be welcome through the project's normal review process, so
that capabilities built for \sysname become available to all ZMap users
instead of living in a fork. He has also offered advice on scanning
methodology, covering rate control, vantage-point selection, probe-module
design, and the safeguards that should govern the observation of operational
networks. This gives Thrust~1 a maintained upstream path and review capacity
the project does not otherwise have.

\subsection{The NGINX community and F5, Inc}
Buu Lam, Director of Community Evangelism for NGINX at F5, operates the
channels through which NGINX practitioners exchange operational knowledge. He
has offered to help the team reach practitioners willing to describe
representative deployment topologies, which is the source of the deployment
reality that Thrust~1 requires; to review configuration-level findings and
assess whether a given result reflects a pattern common in the field or an
artifact of a synthetic test network; and to make NGINX community channels,
events, and technical-content programs available for carrying results and
artifacts to operators. These arrangements are subject to internal review at
F5.

\subsection{The RabbitMQ project and Broadcom Inc}
Gabriele Santomaggio, a member of the RabbitMQ team at Broadcom, has offered
maintainer input on the properties that determine whether a replicated broker
is faithful, including which aspects of cluster formation and inter-node
communication are security-relevant, how virtual-host and permission
structures should be represented, which plugin combinations change the exposed
surface, and where a superficially correct replica would diverge from a real
deployment under load or partition. This supports the clustered-service stress
test of Thrust~2 directly. He has further indicated that findings in the code
would be handled through the project's normal security reporting process, and
that findings about deployment and configuration would inform safer defaults,
configuration validation and diagnostics, and corrections to the documentation
and reference deployments that users copy. These arrangements are subject to
internal review at Broadcom.

\section{Other resources}

\subsection{Competition infrastructure}
The PIs have organized the iCTF competition annually since 2003. The twinCTF
competition of Section~11 will reuse the existing iCTF hosting and scoring
infrastructure rather than standing up new infrastructure of its own.

\subsection{Shared development infrastructure}
The two institutions already share a single source repository and continuous
integration system for existing joint efforts, and this project will use the
same arrangement, as described in Section~6.

\end{document}
